\chapter{Overview and workflow selection}
\label{chap:orientacion}

\section{System components}

The AKD1000 does not operate as a standalone application. The host runs the operating
system and runtime; PCIe carries commands and data; the driver exposes the device to
the system; and the FBZ model defines the network that the runtime can map.
Separating these layers avoids looking for a power fault in Python or reinstalling
the driver to fix a missing package.

\begin{table}[H]
\centering
\caption{Role of each system layer.}
\label{tab:capas}
\begin{tabularx}{\textwidth}{@{}L{0.22\textwidth}YY@{}}
\toprule
\textbf{Layer} & \textbf{Role} & \textbf{What it does not prove on its own} \\
\midrule
Host & Runs Linux or Windows, Python, data preparation and logging. & That the computation ran on the board. \\
PCIe interface & Carries configuration, inputs and outputs. & That a usable driver is available. \\
Linux driver & Provides device access, usually through \code{/dev/akida*}. & That the model is compatible. \\
Akida runtime & Loads FBZ files, discovers devices, maps and executes models. & That virtual mapping is physical inference. \\
AKD1000 & Runs the sequences assigned to the hardware. & Total Python execution time or total platform power consumption. \\
FBZ & Stores a serialized Akida network. & The meaning of inputs, labels or the dataset. \\
\bottomrule
\end{tabularx}
\end{table}

The PCIe board documented by BrainChip contains an AKD1000, uses a PCIe host interface
and has LPDDR4 memory accessible through DMA \cite{brainchip_card_brief}.
The SoC product brief describes twenty neural processing cores and a two-lane PCIe 2.1
endpoint \cite{brainchip_soc_brief}. These are product specifications, not a guarantee
of latency, energy or resource utilization for every network.

\section{Three decisions before you start}

\paragraph{Decision 1: software or hardware.}
Without a board, follow route \textbf{S}: software environment, test model and virtual
device. With a board and a documented PCIe connection, complete route S first, then
continue with route \textbf{H}. Simulation checks the model and package; it does not
check the cabling, driver or board.


\paragraph{Decision 2: current or reference environment.}
To learn with the current documentation, use route \textbf{A}: MetaTF 2.19.3, or the
version listed on the official page when you read this manual. To repeat the included
test, use route \textbf{R}: Python 3.11 and Akida 2.19.1. Do not combine both installations in one virtual environment.


\paragraph{Decision 3: host type.}
A Linux PC with a PCIe slot and a Raspberry Pi 5 with an FFC adapter need different
assemblies. Windows is included in the current MetaTF software configurations
\cite{brainchip_install_2193}, but this manual does not claim physical board support
on Windows: the public PCIe driver instructions consulted are for Linux
\cite{brainchip_driver}.


\section{Workflow matrix}

\begin{longtable}{@{}L{0.12\textwidth}L{0.21\textwidth}L{0.29\textwidth}L{0.28\textwidth}@{}}
\caption{Workflows covered by the manual.}\label{tab:rutas}\\
\toprule
\textbf{Route} & \textbf{Host} & \textbf{Objective} & \textbf{Final evidence} \\
\midrule
\endfirsthead
\toprule
\textbf{Route} & \textbf{Host} & \textbf{Objective} & \textbf{Final evidence} \\
\midrule
\endhead
S & Compatible Windows or Linux host, no board & Install MetaTF, import Akida and run in software. & Recorded version, model output and virtual mapping. \\
H-PC & Linux host, board and documented PCIe connection & Detect the AKD1000 and run on it through the public driver. & \code{lspci}, module, node, \code{akida devices} and physical test. \\
H-Pi & Raspberry Pi 5, FFC adapter and board & Enable PCIe when needed and complete H-PC. & Boot configuration, enumeration and physical test. \\
R & Python 3.11 and reference SDK & Repeat the smoke test with Akida 2.19.1. & Known output and verified virtual mapping. \\
\bottomrule
\end{longtable}

\section{Check each stage}

Each stage answers a different question:

\begin{enumerate}[nosep]
  \item \textbf{Inventory:} do I know the exact host, board, adapter, power supplies,
        operating system, kernel and architecture?
  \item \textbf{Assembly:} is the chain connected according to the documentation for those revisions?
  \item \textbf{Enumeration:} does the host see a new PCIe endpoint?
  \item \textbf{Driver:} has the module been built for this kernel, and is the node accessible?
  \item \textbf{SDK:} does the active interpreter import the expected Akida version?
  \item \textbf{Model:} does the FBZ load, have the expected contract and map as intended?
  \item \textbf{Execution:} does a known input produce a known output on hardware?
  \item \textbf{Measurement:} are the window, units and power measurement boundary stated?
\end{enumerate}

Resolve the first failed stage before continuing. Save its error to help locate
the fault.
